%=============================================================================!
% 6.A  ANSWERS TO THE CHECKS                                                  !
%=============================================================================!
\unnumbered{Answers to the checks}
\label{sec:la-answers}

\answerto{sec:la-coordinates}The double pendulum in a plane needs the two
angles of its rods, 2 coordinates. The rigid diatomic molecule needs the
three coordinates of its centre of mass and two angles for the direction
of its bond, 5, the $3N - 1$ left when one distance of $N = 2$ atoms is
fixed.

\answerto{sec:la-action}With $U = 0$ the expansion of the action has only
the kinetic terms: $\Act(\zeta) = \Act(0) + \zeta\int m\dot{x}\dot\eta\,
\dd t + \zeta^2\,\frac12 m\int\dot\eta^2\,\dd t$. On the straight path
$\dot{x}$ is constant, so the middle term is $m\dot{x}\,[\eta]_{t_1}^{t_2}
= 0$, the bump being zero at both ends. The last term is positive unless
$\dot\eta = 0$ throughout, which with $\eta$ zero at the ends means no
change at all. Any other path with the same ends is the straight one plus such a change,
with $\zeta = 1$ and $\eta$ the difference, and with $U = 0$ the expansion
is exact, so every other path has more action.

\answerto{sec:la-euler}$\partial\Lag/\partial\dot{x} = m\dot{x}$ and
$\partial\Lag/\partial x = -kx$, so $m\ddot{x} = -kx$, the equation of
\cref{sec:ne-spring}.

\answerto{sec:la-constrained}From \SI{0.554}{\second} to
\SI{1.533}{\second}, that is $1.533 - 0.554 = \SI{0.979}{\second}$.

\answerto{sec:la-polar}$\Lag$ does not contain $x$, only $\dot{x}$, so
$p_x = \partial\Lag/\partial\dot{x} = m\dot{x}$ is conserved: nothing
pushes the ball sideways. It does contain $y$, through $mgy$, so $p_y =
m\dot{y}$ changes, at the rate $\partial\Lag/\partial y = -mg$, the weight.

\answerto{sec:la-multipliers}With $\theta_0 = \ang{90}$, $\cos\theta_0 =
0$, and at the bottom $\cos\theta = 1$, so $F_{\mathrm{rod}} = mg(3 - 0) =
3mg$, three times the weight.
